% R32 section 5: periodic boundary conditions (input by R32_floquet_spectra.tex)
% Numbers come from the generated tables tab_r32_bc_H_N14.tex and tab_r32_shape_H_pbc_N14.tex
% (qca_fragmentation.spectral.analysis.write_bc_table / write_shape_table).
The user asked whether the classes survive a change of boundary condition. The obvious suspects are the one-sided
rules, whose block tower is anchored at the open ends. All sixteen rules were run on rings of even length $N=6\ldots14$. At $N=14$, the eight rules with $k\ge2$ that are not
arithmetic have one $\approx16\,000$-dimensional sector each, and each run took about 16~min.

\subsection{What is resolved on the ring, and which ensemble to expect}
On an even ring, translation by one site $T_1$ maps layer $A$ onto layer $B$, for every rule. The half-step
\begin{equation}
  \Omega_T = T_1 L_A,\qquad [\Omega_T,U]=0,\qquad \Omega_T^2 = U\,T_2,
\end{equation}
is therefore a symmetry of every rule. Its eigenphase $\varphi$ fixes both the quasienergy and the momentum $j$ of the
two-site translation $T_2$ ($e^{2i\varphi}=e^{-i\theta}e^{2\pi i j/(N/2)}$). A statistics block is
(Krylov sector, gate-algebra block, $j$), and its level variable is $\varphi$. The site reflection $R$ does \emph{not}
commute with $\Omega_T$ and is not resolved. Which random-matrix ensemble applies depends on the type of sector:
\begin{itemize}
\item \emph{Real sectors} ($j=0$, and $j=N/4$ when $N/4$ is an integer) are mapped to themselves by complex conjugation
  $K$. At $j=0$, $T_1$ is an involution, so $\Omega_T$ is again a product of two reflections: $\beta=1$ on the half
  spectrum (rule~57: $0.549$ at $N=14$). At $j=N/4$, $T_1^2=-1$ and the reflection structure is lost, so the reference is
  a real orthogonal matrix without that structure, which has $\beta=2$-like local statistics (rule~57: $0.610$ at
  $N=12$). For mirror-symmetric rules, the real sectors also contain the unresolved $\mathbb Z_2$ generated by $L_AR$,
  which pushes $\langle r\rangle$ down. \textbf{The real sectors are therefore not used as the diagnostic.}
\item \emph{Complex sectors} ($j\ne -j$) carry most of the levels, so they are the diagnostic. $K$ maps $j\to-j$, so a
  time-reversal symmetry inside sector $j$ needs a unitary $P$ that also flips $j$ and satisfies $PUP^{-1}=U^{-1}$. For a
  mirror-symmetric rule the site reflection does this ($R$ exchanges the layers), so the naive reference is COE ($\beta=1$,
  $0.5307$). For a rule without mirror symmetry the naive reference is CUE ($\beta=2$, $0.5996$).
\end{itemize}
The block classification (real or complex) is numerical: a block counts as real if its eigenphases are closed under
conjugation to $10^{-7}$.

\paragraph{Correction during the analysis.} The first version of this conjugation test compared the sorted lists of $x$
and $-x$ element by element. A level at $0\equiv2\pi$ (the $j=0$ sector has forced eigenvalues $\pm1$) moves to the
other end of the sorted list and shifts every pair. The test then wrongly classified $j=0$ as complex, so its full circle
was used. The fix (\fn{stats.conj\_symmetric}) matches each $-x$ to its nearest level on the circle. A test pins the
$j=0$ block of rule~57 at $N=8$ against the exact Python engine. The obc0 numbers of Sec.~\ref{sec:results} were
recomputed: none of the 145 units changed.

\subsection{Results}
\begin{table}[h]
\centering\small
\resizebox{\textwidth}{!}{\begin{tabular}{llrrrrrrrlrr}
\toprule
 & & \multicolumn{2}{c}{sectors} & \multicolumn{2}{c}{blocks} & obc & \multicolumn{3}{c}{pbc $\langle r\rangle$} & \multicolumn{2}{c}{singlets} \\
$k$ & $W$ & obc & pbc & obc & pbc & $\langle r\rangle_\mathrm{full}$ & real $j$ & complex $j$ & naive ref. & obc & pbc \\
\midrule
0 & 204 & 16384 & 16384 & 16384 & 16384 & -- & -- & -- & COE & 0 & 0 \\
1 & 108 & 3329 & 2627 & 5694 & 4517 & 0.383 & 0.404 & 0.386 & COE & 2070 & 1667 \\
1 & 156 & 987 & 844 & 5740 & 5072 & -- & -- & -- & CUE & 2439 & 2016 \\
1 & 198 & 987 & 844 & 5740 & 5072 & -- & -- & -- & CUE & 2439 & 2016 \\
1 & 201 & 1897 & 2627 & 4143 & 4517 & 0.395 & 0.403 & 0.452 & COE & 1760 & 1667 \\
2 one-sided & 60 & 15 & 2 & 105 & 9 & 0.395 & 0.414 & 0.511 & CUE & 24 & 1 \\
2 one-sided & 102 & 15 & 2 & 105 & 9 & 0.395 & 0.414 & 0.511 & CUE & 24 & 1 \\
2 one-sided & 153 & 1 & 2 & 15 & 9 & 0.421 & 0.428 & 0.531 & CUE & 2 & 1 \\
2 one-sided & 195 & 1 & 2 & 15 & 9 & 0.421 & 0.427 & 0.531 & CUE & 2 & 1 \\
2 XOR & 105 & 8 & 7 & 15 & 20 & 0.356 & 0.405 & 0.391 & COE & 7 & 7 \\
2 XOR & 150 & 8 & 10 & 15 & 52 & 0.419 & 0.424 & 0.440 & COE & 7 & 6 \\
3 & 54 & 2 & 2 & 3 & 9 & 0.523 & 0.383 & 0.483 & COE & 1 & 1 \\
3 & 57 & 1 & 1 & 3 & 8 & 0.531 & 0.549 & 0.602 & CUE & 2 & 1 \\
3 & 99 & 1 & 1 & 3 & 8 & 0.531 & 0.549 & 0.602 & CUE & 2 & 1 \\
3 & 147 & 1 & 2 & 5 & 9 & 0.527 & 0.424 & 0.518 & COE & 4 & 1 \\
4 & 51 & 1 & 1 & 16384 & 16384 & -- & -- & -- & COE & 16384 & 16384 \\
\bottomrule
\end{tabular}
}
\caption{Open chain against ring at $N=14$, all sixteen rules. ``naive ref.''\ is
the expectation for the complex sectors from the identified symmetries (mirror $\Rightarrow$ COE). The one-sided rules
violate it (Table~\ref{tab:shape}). The $N=12$ table is \fn{tab\_r32\_bc\_H\_N12.tex}.}
\label{tab:bc}
\end{table}

\begin{table}[h]
\centering\small
\begin{tabular}{llrrrrrrrr}
\toprule
$k$ & $W$ & mirror & \#$r$ & $\langle r\rangle$ & $P(r<0.15)$ & $\chi^2_\mathrm{P}$ & $\chi^2_{\beta=1}$ & $\chi^2_{\beta=2}$ & $\chi^2_\mathrm{P+\beta=2}$ \\
\midrule
3 & 57 & no & 13944 & 0.602 & 0.023 & 8482 & 1083 & 47 & 47 \\
3 & 99 & no & 13944 & 0.602 & 0.023 & 8482 & 1083 & 47 & 47 \\
3 & 147 & yes & 13938 & 0.518 & 0.082 & 3444 & 124 & 4324 & 400 \\
3 & 54 & yes & 13938 & 0.483 & 0.105 & 2119 & 850 & 9869 & 466 \\
2 one-sided & 60 & no & 13938 & 0.511 & 0.079 & 3228 & 163 & 4873 & 399 \\
2 one-sided & 102 & no & 13938 & 0.511 & 0.079 & 3228 & 163 & 4873 & 399 \\
2 one-sided & 153 & no & 13938 & 0.531 & 0.063 & 4191 & 27 & 2458 & 385 \\
2 one-sided & 195 & no & 13938 & 0.531 & 0.063 & 4191 & 27 & 2458 & 385 \\
\bottomrule
\end{tabular}

\caption{Shape of $P(r)$ in the complex momentum sectors of the ring, $N=14$: $\chi^2$ of a 10-bin histogram (9 degrees
of freedom) against the Atas surmises, and against the best Poisson+CUE mixture. A crossover from Poisson towards CUE
can pass through the COE \emph{mean}, but then its small-$r$ tail is too heavy ($P(r)\sim r^\beta$). The surmise values
of $P(r<0.15)$ are 0.261 (Poisson), 0.064 ($\beta=1$) and 0.019 ($\beta=2$). The surmise deviates from large-$N$ random
matrices by about 1\%, so $\chi^2$ of a few tens at $n\approx1.4\times10^4$ is a good fit.}
\label{tab:shape}
\end{table}

\begin{figure}[h]
\centering
\ifgraphic{../../figures/r32/fig_r32_bc_compare_H.pdf}{\textwidth}
\caption{Open chain against ring, each rule at the largest even $N$ available for both boundaries ($N=14$ for the
$k\ge2$ non-arithmetic rules, $N=12$ otherwise). (a) Resolved $\langle r\rangle$. For the ring, the real momentum sectors
(open squares) and the complex sectors (filled squares) are shown separately. (b) Number of resolved blocks.}
\label{fig:bc}
\end{figure}

\paragraph{$k=3$: the class survives, and the ensemble follows the symmetry.}
57 and 99 have no mirror symmetry and become CUE on the ring: $\langle r\rangle=0.602\pm0.002$, with a GUE-shaped
$P(r)$ ($\chi^2=47$, against 1083 for $\beta=1$). This is the textbook outcome for a chaotic system once its momentum
sectors are complex. 147 (mirror) is COE-shaped and approaches $0.531$ from below ($0.469$, $0.496$, $0.518$ for
$N=10,12,14$). 54 (mirror) follows more slowly ($0.458$, $0.472$, $0.483$). On the open chain, 54 also lagged the other
three rules by about two sizes.

\paragraph{One-sided rules: the tower is a boundary effect, and on the ring they are chaotic.}
On the ring the one-sided rules have 2 Krylov sectors and 9 blocks, against a tower of $N+1$ blocks (60: 105 blocks at
$N=14$). In the complex sectors their statistics are $\beta=1$ in \emph{shape}:
\begin{itemize}
\item 153 and 195: $\langle r\rangle=0.531$, $P(r<0.15)=0.063$, and $\chi^2_{\beta=1}=27$ against 385 for the best
  Poisson+CUE mixture.
\item 60 and 102: still rising ($0.423$, $0.471$, $0.511$), and already closer to COE than to any mixture.
\end{itemize}
This contradicts the naive reference: these rules are \emph{not} mirror symmetric, so one expects CUE. There must
therefore be an antiunitary symmetry inside each momentum sector that we have not identified. A search at $N=8$ found nothing beyond the known layer swap $T_1UT_1^{-1}=U^{-1}$, which preserves $j$ and so gives no
time reversal inside a sector. It looked for $W$ with $WUW^{-1}=U^{\pm1}$, and separately with
$W\Omega_TW^{-1}=\Omega_T^{-1}$, over $W=D\,L\,R^aT_s$: $a\in\{0,1\}$, all $s$, $L$ the identity or a uniform or
sublattice product of $X$ or $H$ (and of $Z$ without $D$), and $D$ an \emph{arbitrary} diagonal phase operator, solved
exactly as a consistency problem on the support graph. For 147 the search returns the reflection, which validates it.
The symmetry is an open point
(Sec.~\ref{sec:open-pbc}).

The consequence for the classification is that \textbf{the one-sided class is defined by its boundary}. With open ends
the tower splits the space into $N+1$ blocks and the statistics stay near Poisson. On the ring the same bulk dynamics is
chaotic.

\paragraph{$k=2$ XOR: still arithmetic.} The ring makes the spectra of 150 and 105 \emph{more} degenerate. At $N=14$,
150 has 255 distinct quasienergies and 105 has 251, against $2187=3^7$ for both on the open chain. Their statistics
are not those of a chaotic spectrum either: $\langle r\rangle=0.440$ (150) and $0.391$ (105) in the complex sectors,
with a small-$r$ tail of $P(r<0.15)=0.22$ and $0.45$ (Poisson: 0.26; $\beta=1$: 0.06). That tail is the signature of
near-coincident levels in an arithmetic spectrum (Table~\ref{tab:shape}). The spectra are arithmetic, not chaotic, as on
the open chain.

\paragraph{$k=1$: still exponentially fragmented.} At $N=14$ the ring has 2627 sectors for 108 and 201 (open chain:
3329 and 1897) and 844 for 156 and 198 (open: 987). The blocks stay small: the largest has 143 states for 108/201 and 25
for 156/198. On the ring, 108 and 201 have \emph{identical} block structures (sectors, blocks, singlets and the count
of distinct levels), but different statistics ($0.399$ and $0.417$, from 2547 ratios each). This is consistent with a
relabelling of the basis that maps one rule's sector graph onto the other's, while the gate $H$ is not invariant under
it. The same holds for 156 and 198, which are mirror images. Within the large uncertainty of such small blocks, the
statistics stay between Poisson and the structured value of the open chain.

\paragraph{Exact singlets on the ring.} Every $k=3$ and one-sided rule has exactly \emph{one} gate singlet on the ring, at
$\theta=0$, for every $N$ (the XOR rules keep about one per sector, as on the open chain). The singlets at $\theta=\pi$ of the open chain are gone, so they are boundary states. For
57, 99, 147, 153 and 195 the ring singlet is a product state ($S_\mathrm{vN}=0$). For 54, and also for 60 and 102, it is
entangled, with the \emph{same} entropy as the entangled singlet of 54 on the open chain:
$S_\mathrm{vN}=0.4045,\ 0.3447,\ 0.2908,\ 0.2430$ for $N=8,10,12,14$. Equal entropies across two rules and two boundary
conditions suggest one and the same state. The eigenvectors were not stored, so overlaps could not be checked.

\subsection{What survives}
\label{sec:open-pbc}
\begin{center}
\begin{tabular}{lp{0.3\textwidth}p{0.42\textwidth}}
\toprule
class & open chain & ring \\
\midrule
$k=3$ & chaotic (COE) & chaotic: CUE (57, 99) or COE (147, 54), as the symmetry dictates \\
$k=2$ one-sided & boundary tower, near Poisson & tower gone, chaotic ($\beta=1$ shape) \\
$k=2$ XOR & arithmetic & arithmetic, even more degenerate \\
$k=1$ & exponential fragmentation & unchanged \\
$k=0,4$ & trivial & trivial \\
\bottomrule
\end{tabular}
\end{center}
The $k=3$, XOR, $k=1$ and trivial classes are robust. The one-sided class is not a bulk class: its non-ergodicity comes
from the open ends. ``Chaos needs three firing patterns'' is therefore a statement about the open chain. On the ring,
two firing patterns suffice if they depend on a single neighbour.

\paragraph{Open points.}
\begin{enumerate}
\item The antiunitary symmetry behind the $\beta=1$ statistics of the one-sided rules on the ring.
\item Resolving $L_AR$ in the real sectors of the mirror rules, which needs the stored eigenvectors or one extra
  symmetry in the Julia pipeline.
\item The identity of the entangled singlet shared by 54 (both boundaries) and 60/102 (ring).
\end{enumerate}
